\chapter{KẾT LUẬN VÀ KIẾN NGHỊ}
\label{ch:conclusion}

Dự án \textbf{Tech-Green} là minh chứng cho việc kết hợp hài hòa giữa tư duy kinh doanh thực tế và năng lực công nghệ chuyên sâu của sinh viên thời đại số. Không chỉ dừng lại ở ý tưởng, dự án đã xây dựng được một khung vận hành hoàn chỉnh, sẵn sàng triển khai thực tế.

\section{Tổng kết các kết quả đạt được}

\subsection{Về mặt Kinh doanh và Vận hành}
\begin{itemize}
    \item \textbf{Xác định đúng thị trường ngách:} Tech-Green đã lựa chọn phân khúc sản phẩm Decor bàn làm việc (Desk Setup) đang tăng trưởng nóng, giải quyết đúng "nỗi đau" của khách hàng về không gian sống.
    \item \textbf{Mô hình tinh gọn:} Áp dụng triệt để nguyên lý \textit{Lean Startup}, tối ưu hóa quy trình nhập hàng từ 1688 và vận hành kho bãi tại nhà để giảm thiểu chi phí cố định.
\end{itemize}

\subsection{Về mặt Công nghệ và Đổi mới sáng tạo}
Đây là điểm sáng lớn nhất của dự án so với các đối thủ truyền thống:
\begin{itemize}
    \item \textbf{Làm chủ công nghệ lõi:} Tự xây dựng hệ thống E-commerce bằng Python (Django) thay vì phụ thuộc vào nền tảng có sẵn.
    \item \textbf{Cá nhân hóa trải nghiệm:} Ứng dụng thành công thuật toán gợi ý sản phẩm (Content-based Filtering) và Chatbot NLP để tư vấn bán hàng tự động, mang lại trải nghiệm mua sắm vượt trội.
\end{itemize}

\section{Định hướng phát triển trong tương lai (Future Roadmap)}

Tech-Green không chỉ dừng lại ở một cửa hàng bán lẻ, mà hướng tới việc xây dựng một hệ sinh thái \textbf{"Green-Tech"}.

\begin{description}
    \item[Giai đoạn 1 (Năm 2026 - 2027):] Ổn định dòng tiền, đạt điểm hòa vốn sau 6 tháng. Hoàn thiện hệ thống Website và mở rộng kênh TikTok Shop.
    \item[Giai đoạn 2 (Năm 2027 - 2028):] Phát triển ứng dụng di động (Mobile App) sử dụng Flutter/React Native, tích hợp tính năng AR (Thực tế tăng cường) cho phép khách hàng ướm thử cây vào bàn làm việc qua Camera điện thoại.
    \item[Giai đoạn 3 (Tầm nhìn 2030):] Mở rộng sang mảng B2B (Doanh nghiệp), cung cấp giải pháp "Văn phòng xanh" trọn gói cho các công ty công nghệ và Co-working space.
\end{description}

\section{Kiến nghị và Đề xuất}

Để dự án có thể chuyển mình từ đồ án sinh viên thành một doanh nghiệp khởi nghiệp bền vững, nhóm tác giả xin đưa ra một số kiến nghị:

\begin{enumerate}
    \item \textbf{Đối với Nhà trường:} Kính mong Nhà trường hỗ trợ không gian ươm tạo (Co-working space) tại trường để nhóm có địa điểm làm việc và trưng bày sản phẩm mẫu, giúp tiết kiệm chi phí thuê mặt bằng giai đoạn đầu.
    \item \textbf{Đối với các đơn vị hỗ trợ khởi nghiệp:} Hỗ trợ tư vấn pháp lý về Sở hữu trí tuệ (đăng ký nhãn hiệu Tech-Green) và kết nối với các quỹ đầu tư thiên thần (Angel Investors) quan tâm đến lĩnh vực Retail Tech.
\end{enumerate}

\section*{Lời kết}
Tech-Green mang trong mình khát vọng của người trẻ về việc ứng dụng công nghệ để giải quyết các vấn đề cuộc sống. Chúng tôi tin rằng, với sự chuẩn bị kỹ lưỡng về kế hoạch kinh doanh và nền tảng kỹ thuật vững chắc, Tech-Green hoàn toàn có khả năng vươn lên trở thành thương hiệu Top-of-mind về Decor công nghệ tại Việt Nam.